\chapter{Introduzione}
\label{cap:introduzione}

\section{L'azienda ospitante}

\begin{figure}[!h] 
    \centering 
    \includegraphics[width=0.7\columnwidth]{Omicron_Logo} 
    \caption{Logo Omicron}
\end{figure}

\noindent L’azienda Omicron Consulting è specializzata nello sviluppo di software gestionali e di revisione di processi aziendali. Essa è presente nel mercato ICT dal 1980, spiccando su vari settori, in cui hanno effettuato importanti implementazioni in area ICT come:  Manufacturing, Automotive, Aerospace, Logistics e altre aree.
Con particolare riferimento al settore Manufacturing si sono specializzati nello sviluppo di progetti di trasformazione ERP (acquisendo la certificazione VAR di SAP), stringendo alleanze strategiche con realtà ICT nazionali ed internazionali.\\
Offrono servizi di gestione e supporto di sistemi ERP, sviluppo di progetti di Business Intelligence e personalizzazione di sistemi software.\\
Per completare il pacchetto dei servizi offerti, Omicron ha un'esperienza di alto livello negli ambiti Banking, Finance and Insurance, lavorando con le principali istituzioni bancarie e assicurative italiane.\\
Omicron oltre alle offerte che dedica ai clienti, si occupa anche di garantire un'alta formazione delle proprie risorse, investendo su progetti di ricerca e sviluppo.

\section{Il progetto}
\subsection{Presentazione}
L'applicativo permette la creazione di richieste di pianificazioni di risorse aziendali per svolgere un determinato incarico. In questo progetto per risorse aziendali si fa sempre riferimento alle risorse umane, quindi al personale o alla forza lavoro dell'azienda.\\
Il Project Manager può inoltrare richieste per l'allocazione di risorse aziendali per svolgere una task rapida o un incarico relativo ad un progetto più ampio, chiedendo disponibilità di figure professionali con determinate skill richieste ed ulteriori parametri come l'effort e la seniority.\\ 
In una fase successiva, il Program Manager prende visione delle richieste effettuate dal Project Manager e si incarica di creare una Pianificazione per ciascuna figura professionale richiesta. Questa fase prevede l'assegnazione delle risorse adeguate a ciascuna richiesta, specificando parametri quali la durata, l'attività da svolgere e molti altri aspetti rilevanti per garantire una gestione efficiente delle risorse umane aziendali.\\

\subsection{Motivazione del progetto}
%L'approccio adottato per la gestione delle pianificazioni delle risorse e la loro disponibilità veniva gestito attraverso fogli Excel compilati e rivisti dai Program manager. Le nuove richieste di pianificazioni vengono comunicate ai Program Manager tramite posta elettronica, telefono e chat, rendendo arduo tenere traccia di tutto.\\
%Il progetto nasce dunque dall’esigenza di semplificare la gestione delle richieste e delle pianificazioni, garantendo una visione più rapida della disponibilità delle risorse.\\

L'approccio tradizionale adottato per la gestione delle pianificazioni delle risorse e la loro disponibilità all'interno dell'azienda era basato su fogli Excel compilati e costantemente rivisti dai Program Manager. Le nuove richieste di pianificazioni provenienti dai Project Manager venivano comunicate ai Program Manager attraverso diversi canali, tra cui posta elettronica, telefono e chat. Questo metodo di comunicazione rendeva spesso arduo il compito di tenere traccia di tutte le richieste e di coordinare efficacemente le risorse umane.\\
Il progetto nasce dunque dall’esigenza di semplificare e modernizzare la gestione delle richieste e delle pianificazioni delle risorse aziendali. L'obiettivo principale è garantire una visione più rapida, accurata e centralizzata della disponibilità delle risorse, riducendo al minimo i potenziali conflitti e le sovrapposizioni nelle pianificazioni.\\

\subsection{Il mio ruolo nel progetto}
%L'obiettivo di questa attività di stage era quello di inserirmi all'interno del progetto, più precisamente nello sviluppo del lato back-end sviluppando un'API REST e le tabelle del database su cui poi l'API poggiava. Le tabelle dovevano inoltre integrarsi con le tabelle del database aziendale, recuperando dati utili alle funzionalità.\\
%Il mio lavoro consisteva nel sviluppare la possibilità di effettuare operazioni CRUD tramite gli endpoint dell'API su: richieste di figure professionali, pianificazioni e milestone, oltre a operazioni di sola lettura sulle anagrafiche dell'azienda.\\
%Dopo aver studiato le tecnologie necessarie e svolto alcune esercitazioni, ho iniziato lo sviluppo del prodotto basandomi su due documenti: un documento mock\textsubscript{g} dell'API e un'Analisi dei Requisiti che descriveva il prodotto nella totalità.

L'obiettivo principale di questa attività di stage era quello di inserirmi all'interno del progetto aziendale, concentrandomi specificamente sullo sviluppo del lato back-end. Il mio compito principale consisteva nell'implementare un'API REST e le relative tabelle del database su cui l'API si basava. Queste tabelle dovevano essere integrate con il database aziendale esistente, permettendo il recupero dei dati necessari per il funzionamento delle varie funzionalità.\\
Il mio lavoro consisteva nel sviluppare le funzionalità che consentivano operazioni di Creazione, Lettura, Aggiornamento e Cancellazione (CRUD) relative a richieste di figure professionali, pianificazioni e milestone da associare alle pianificazioni, oltre a operazioni di sola lettura sulle anagrafiche dell'azienda.\\
Prima di immergermi nell'implementazione, ho dedicato tempo a studiare le tecnologie e gli strumenti necessari per il progetto completando esercitazioni mirate per acquisire familiarità con le nuove tecnologie e approcci di sviluppo.\\
Durante la fase di sviluppo, mi sono basato su due documenti forniti: un documento mock\textsubscript{g} dell'API e un'Analisi dei Requisiti che descriveva il prodotto nella sua totalità.

%Le funzionalità da me sviluppate principalmente sono due: le richieste e le pianificazioni di risorse aziendali.\\
%Questo è stato realizzato attraverso la creazione di nuove tabelle da inserire nel database aziendale per supportare l'utilizzo dell'applicativo e lo sviluppo dell'API\textsubscript{g} REST. Gli endpoint dell'API consentono operazioni CRUD\textsubscript{g} su richieste, pianificazioni e milestone, oltre a operazioni di sola lettura sulle anagrafiche dell'azienda.\\
%Dopo aver studiato le tecnologie necessarie, ho iniziato lo sviluppo del prodotto basandomi su due documenti: un documento mock\textsubscript{g} dell'API e un'Analisi dei Requisiti che descriveva il prodotto finale.


\section{Organizzazione del testo}
Questa sezione è dedicata alla spiegazione della struttura del documento, per dare indicazioni su come è organizzato il testo.
\begin{description}

\item [{\hyperref[cap:introduzione]{Capitolo 1: Introduzione}},] introduce il progetto, il mio ruolo all'interno del progetto e il profilo aziendale;

    \item [{\hyperref[cap:analisi-requisiti]{Capitolo 2: Analisi dei Requisiti}},] descrive i casi d'uso individuati ed i relativi requisiti;
    
    \item[{\hyperref[cap:tecnologie]{Capitolo 3: Background Tecnologico}},] descrive le tecnologie studiate e gli strumenti utilizzati per sviluppare il prodotto;
    
    \item[{\hyperref[cap:progettazione]{Capitolo 4: Progettazione}},] descrive il periodo di progettazione dell'architettura su cui poggia il sistema e i concetti fondamentali;
    
    \item[{\hyperref[cap:database]{Capitolo 5: Modello dati}},] parla della progettazione del database e delle tabelle utilizzate spiegandone l'utilizzo;
    
    \item[{\hyperref[cap:api-rest]{Capitolo 6: Sviluppo API REST}},] tratta dell'API sviluppata e della sua implementazione;
    
    \item[{\hyperref[cap:codifica]{Capitolo 7: Codifica}},] descrive la codifica dell'API REST, l'organizzazione del codice, i moduli sviluppati e le parti più rilevanti;
    
    \item[{\hyperref[cap:verifica-validazione]{Capitolo 8: Verifica e Validazione}},] descrive le fasi di Verifica e Validazione del prodotto;
    
    \item[{\hyperref[cap:deploy]{Capitolo 9: Deploy dell'applicazione}},] descrive la fase di Deploy dell'applicazione e di come è stata svolta col tutor;
    
    \item[{\hyperref[cap:conclusioni]{Capitolo 10: Conclusioni}},] tratta delle considerazioni personali sull'attività svolta;
    
\end{description}

\section{Convenzioni tipografiche nel documento}
Durante la stesura del testo sono state adottate le seguenti convenzioni:
\begin{itemize}
\item la sezione del glossario contiene i termini ritenuti ambigui o non di uso comune, che necessitano quindi di una loro definizione. Il suo scopo è quello di fornire una comprensione comune del linguaggio utilizzato e di evitare confusione o interpretazioni errate;
\item per la prima occorrenza di un termine inserito nel glossario viene utilizzata la seguente nomenclatura: termine\textsubscript{g}.\\
\end{itemize}

